\chapter{Integral cyclic comparison and Hamiltonian homotopies}
\label{mht:ch:integral-cyclic}

The coefficient Weyl algebra has cyclic operations that ordinary
traces do not determine. We construct a comparison from its
Hochschild chains to differential forms which preserves the circle
operator. Two independent parameters remain in the comparison.
Their relation connects the vacuum to the volume form before
either parameter is inverted.

Hamiltonian equivariance requires two explicit homotopies.
After passing to continuous duals, the same comparison gives a
criterion for a current to be strictly cyclic. Point currents
then give nonzero scalar classes with separate primitives for
the two parameters.

\section{Two operators on relative chains}

Let $A=\mathbb C[q,p]$ with the direct-limit topology of its
finite-degree subspaces. The second coefficient domain will be
$A=\mathbb C[[q,p]]$ with its formal topology. In either case set
$R_\epsilon=\mathbb C[[\epsilon]]$ and
\[
 A_\epsilon=(A[[\epsilon]],\star),\qquad
 a\star b=m\exp(\epsilon P/2)(a\otimes b),\qquad
 P=\partial_q\otimes\partial_p-\partial_p\otimes\partial_q.
\]
Here $m$ is ordinary multiplication and $|\epsilon|=0$.
This is a deformation of the coefficient algebra, with
$q\star p-p\star q=\epsilon$. It is distinct from the matrix-entry
quantization in Chapter~\ref{mht:ch:matrix-branes}.

Parameter series are completed coefficientwise. For the
polynomial domain, every parameter coefficient is a polynomial,
with no uniform degree bound. For the formal domain, tensor
products are also completed in the separate coordinates of
all slots.

At each parameter order the Moyal product and the circle comparison
have finite differential order in their coefficient slots.
They are continuous on both domains.
Continuity of the chosen cone contractions, coefficientwise
antiderivatives and flow operators is justified separately below.
No finite-differential-order assertion is made for them.

Use normalized Hochschild chains relative to $R_\epsilon$:
\[
 C_n=A_\epsilon\widehat\otimes_{R_\epsilon}
          (A_\epsilon/R_\epsilon1)^{\widehat\otimes n}.
\]
The chain $a_0[a_1|\cdots|a_n]$ has cohomological degree $-n$.
A scalar in a bar slot is zero. The Hochschild boundary is
\begin{align*}
 b(a_0[a_1|\cdots|a_n])
 ={}&(a_0\star a_1)[a_2|\cdots|a_n]\\
 &+\sum_{i=1}^{n-1}(-1)^i
       a_0[a_1|\cdots|a_i\star a_{i+1}|\cdots|a_n]\\
 &+(-1)^n(a_n\star a_0)[a_1|\cdots|a_{n-1}].
\end{align*}
It vanishes on $C_0$.

The Connes operator is
\[
 B(a_0[a_1|\cdots|a_n])
 =\sum_{i=0}^n(-1)^{ni}
      1[a_i|\cdots|a_n|a_0|\cdots|a_{i-1}].
\]
Associativity cancels the pairs of terms in $b^2$.
Cyclic rotation and normalization give $B^2=bB+Bb=0$.
Thus $|b|=1$ and $|B|=-1$ define a mixed complex.
Adjoin $u$ of degree two and write
\[
 D_C=b+uB
 \quad\text{on}\quad C_\bullet[[u]].
\]
This is the complete negative-cyclic chain complex.

Place a differential $n$-form in degree $-n$. Put
$\pi=\partial_q\wedge\partial_p$, use
$\iota_\pi(dq\wedge dp)=1$, and define
$\delta=\iota_\pi d-d\iota_\pi$. In coordinates,
\[
 \delta(a\,dq+b\,dp)=b_q-a_p,\qquad
 \delta(c\,dq\wedge dp)=-dc,\qquad \delta(f)=0.
\]
Direct differentiation gives
$d^2=\delta^2=d\delta+\delta d=0$. The form complex is
\[
 \bigl(\Omega_A^\bullet[[\epsilon,u]],D_\Omega\bigr),
 \qquad D_\Omega=\epsilon\delta+ud.
\]
Write $\Omega=dq\wedge dp$.

The smallest mixed relation already contains both parameters.
With $\chi=1[q|p]-1[p|q]$, one has
\begin{equation}\label{mht:eq:cyclic-vacuum-volume}
 \begin{aligned}
 D_C(q[p])&=\epsilon1+u\chi,\qquad &b\chi=B\chi&=0,\\
 D_\Omega(q\,dp)&=\epsilon+u\Omega, &\delta\Omega=d\Omega&=0.
 \end{aligned}
\end{equation}
The Hochschild boundary of $\chi$ is a scalar in a bar slot
and therefore vanishes. A comparison preserving the actual
circle operator must carry the first relation to the second.

\section{Integration over a circle}

For $n\geq1$, let
\[
 \Delta_n=\{0<t_1<\cdots<t_n<1\},\qquad t_0=0,
\]
with orientation $dt_1\wedge\cdots\wedge dt_n$.
Let $P_{ij}$ apply $P$ to tensor slots $i,j$.
Its derivatives act on form coefficients and leave their
exterior factors fixed. Define
\begin{equation}\label{mht:eq:circle-comparison}
 \mathcal U_n(a_0[a_1|\cdots|a_n])
 =m\int_{\Delta_n}
   \exp\!\left(\epsilon\sum_{i<j}
                 (\tfrac12+t_i-t_j)P_{ij}\right)
       a_0\,da_1\cdots da_n\,dt_1\cdots dt_n.
\end{equation}

In this formula $m$ identifies the coordinate slots and
multiplies their exterior factors. Set $\mathcal U_0$ equal
to the identity. Exterior degree makes $\mathcal U_n=0$ for
$n>2$. A scalar bar input has zero differential, so the formula
descends to normalized chains.

For polynomial inputs the exponential terminates.
The required moments are rational:
\[
 \int_{\Delta_n}t_1^{a_1}\cdots t_n^{a_n}\,dt
   =\prod_{j=1}^n(a_1+\cdots+a_j+j)^{-1}.
\]
Successive integration proves this formula. In particular,
\[
 \mathcal U_n\bmod\epsilon
     =\frac1{n!}a_0\,da_1\wedge\cdots\wedge da_n.
\]
The volume of the simplex supplies the normalized
Hochschild--Kostant--Rosenberg map.

At one bar slot the operator has a simpler expression:
\begin{equation}\label{mht:eq:circle-onebar}
 \mathcal U_1(a_0[a_1])
 =mS(P)(a_0\,da_1),\qquad
 S(P)=\frac{2\sinh(\epsilon P/2)}{\epsilon P}
 =\sum_{j\geq0}\frac{\epsilon^{2j}P^{2j}}
                         {2^{2j}(2j+1)!}.
\end{equation}
The quotient denotes its power series, whose constant
coefficient is one. For example,
\[
 \mathcal U_1(q^3[p^3])
   =(3q^3p^2+\tfrac32\epsilon^2q)\,dp.
\]
Applying $\epsilon\delta$ gives
$9\epsilon q^2p^2+3\epsilon^3/2$, the Moyal commutator of
$q^3$ and $p^3$. The classical term alone would miss the
cubic correction.

The ordered-circle kernel is the one appearing in
\href{https://arxiv.org/abs/math/0311303v2}
{Feigin--Felder--Shoikhet, Section~2.3}.
Its compatibility with $B$ is the identity of
\href{https://arxiv.org/abs/0804.3887v1}
{Willwacher, Theorem~3}.
We prove the flat formulas with the normalization above.

\section{The mixed comparison}

\begin{theorem}\label{mht:thm:integral-cyclic-comparison}
On both specified coefficient domains,
\[
 \mathcal U b=\epsilon\delta\mathcal U,\qquad
 \mathcal U B=d\mathcal U.
\]
Consequently $\mathcal U$ is a quasi-isomorphism
\[
 (C_\bullet[[u]],b+uB)
 \longrightarrow
 (\Omega_A^\bullet[[\epsilon,u]],\epsilon\delta+ud).
\]
It sends $\chi$ to $\Omega$ and preserves
\eqref{mht:eq:cyclic-vacuum-volume} before either parameter
is inverted.
\end{theorem}
\begin{proof}
At fixed parameter order all operators have constant
coefficients and finite differential order. It therefore
suffices to compare their symbols on
$a_i=\exp(k_i\cdot(q,p))$.
The coefficients in the momenta $k_i$ recover the polynomial
identities and their continuous formal extensions.

Set $K=\sum_i k_i$, $c_{ij}=\pi(k_i,k_j)$,
$\Lambda_i=k_i\cdot(dq,dp)$, and
\[
 S(t)=\sum_{i<j}(\tfrac12+t_i-t_j)c_{ij}.
\]
The integrand of \eqref{mht:eq:circle-comparison} is
$e^{K\cdot(q,p)}e^{\epsilon S(t)}
\Lambda_1\wedge\cdots\wedge\Lambda_n$.
The two elementary identities are
\begin{align*}
 \partial_{t_i}S&=\pi(k_i,K),\\
 \delta(e^{K\cdot(q,p)}\Lambda_1\cdots\Lambda_n)
 &=e^{K\cdot(q,p)}\sum_{i=1}^n(-1)^i\pi(k_i,K)
         \Lambda_1\cdots\widehat\Lambda_i\cdots\Lambda_n.
\end{align*}
Integrating the first identity by parts gives the boundary
faces of the ordered simplex.

At $t_1=0$, the factor $e^{\epsilon c_{01}/2}$ is the first
Moyal multiplication. At $t_i=t_{i+1}$, the two omitted-form
terms combine their differentials to
$\Lambda_i+\Lambda_{i+1}$. At $t_n=1$, the cyclic
multiplication has factor $e^{-\epsilon c_{0n}/2}$ and sign
$(-1)^n$. These are exactly the three types of terms in $b$.

The face sum also proves the equation for $n=3$, where the
interior three-form vanishes. For larger $n$ both sides vanish
by form degree. This proves $\mathcal U b=\epsilon\delta\mathcal U$.

For the second identity, move the marked origin around the
circle. Its possible positions are partitioned by the $n+1$
intervals between labeled inputs. Taking each interval as
the first gives the corresponding cyclic rotation in $B$,
with sign $(-1)^{ni}$ from exterior reordering.
The marked origin has total interval length one.

The resulting form factor is
\[
 (\Lambda_0+\cdots+\Lambda_n)
       \wedge\Lambda_1\wedge\cdots\wedge\Lambda_n.
\]
This is the exterior derivative of the exponential coefficient.
The jump of $\tfrac12+t_i-t_j$ at the origin is precisely
the cyclic reordering of the kernel. Thus
$\mathcal U B=d\mathcal U$, with no additional factorial.

At $\epsilon=0$, the diagonal Koszul resolution on
$q_L-q_R,p_L-p_R$ computes Hochschild homology as differential
forms. Its comparison with the normalized bar resolution gives
the displayed Hochschild--Kostant--Rosenberg projection.
The completed diagonal resolution gives the same assertion
for formal functions and completed chains.

Give each coordinate and its differential weight one.
At fixed weight the normalized undeformed comparison cone is
a finite-dimensional acyclic complex: a positive bar slot has
positive weight, which bounds its length. Choose a contraction
at each weight. Their direct sum preserves polynomial inputs.
Their product is continuous on formal inputs.

The deformed differential differs by terms in $(\epsilon,u)$.
Perturbing this contraction gives a convergent geometric series
in the parameter-adic topology, with continuous coefficients.
It contracts the completed comparison cone and proves the
quasi-isomorphism on both stated domains.
Finally, the two linear bar inputs have no nonzero contraction,
so direct antisymmetrization gives $\mathcal U\chi=\Omega$.
\end{proof}

\section{Cohomology before inversion}

The integral cohomology records information hidden by either
parameter specialization. Write $R=\mathbb C[[\epsilon,u]]$.
For this calculation group the completed even and odd form
degrees. The original cohomological grading is retained by
assigning $|u|=2$ and $|\Omega|=-2$.

\begin{proposition}\label{mht:prop:integral-form-cohomology}
The form complex has
\[
 H^{\mathrm{even}}
   =\frac{R1\oplus R\Omega}{R(\epsilon1+u\Omega)},
 \qquad H^{\mathrm{odd}}=dA,
 \qquad (\epsilon,u)H^{\mathrm{odd}}=0.
\]
Both the vacuum and volume classes are nonzero.
\end{proposition}
\begin{proof}
An even element $f+g\Omega$ is closed exactly when
$d(uf-\epsilon g)=0$.
Subtract the constant spatial coefficients of $f$ and $g$,
leaving $f',g'$. Then $uf'=\epsilon g'$.
Reduction modulo $\epsilon$ shows that $f'$ is divisible by
$\epsilon$. Since neither parameter is a zero divisor,
\[
 f'=\epsilon h,\qquad g'=uh.
\]

Every coefficient of $h$ is a curl $b_q-a_p$.
For example, integrate it with respect to $q$ and take $a=0$.
This integration is continuous on the formal domain.
Thus $f'+g'\Omega=D_\Omega(a\,dq+b\,dp)$.

The remaining constant pair has the relation
$\epsilon1+u\Omega$, obtained from $q\,dp$.
Its quotient embeds in $R$ by
\[
 [a+b\Omega]\longmapsto ua-\epsilon b.
\]
The kernel before taking the quotient is exactly the submodule
$R\cdot(\epsilon,u)\subset R^2$, by the same regular-sequence argument.
The two classes map to $u$ and $-\epsilon$, so neither vanishes.

An odd form $\alpha$ is closed exactly when its curl vanishes.
Indeed the scalar and volume components of $D_\Omega\alpha$
are its curl multiplied by $\epsilon$ and $u$.
Closed one-forms on the polynomial or formal plane are exact.
The odd boundaries are
$u\,df-\epsilon\,dg$, precisely
$(\epsilon,u)dA[[\epsilon,u]]$.
\end{proof}

Specializing this answer need not compute the cohomology of
the specialized complex. At $u=0$, the differential is
$\epsilon\delta$, and its degree-zero cohomology is
$A[[\epsilon]]/\epsilon A[[\epsilon]]$.
In particular $q$ survives there. It does not extend to an
even class with nonzero $u$, because $D_\Omega q=u\,dq$.
The vacuum--volume relation alone does not account for this
failure of ordinary base change.

\section{Hamiltonian operations}

Let $\mathfrak h_\epsilon=A[[\epsilon]]/R_\epsilon1$ with
bracket
\[
 \{f,g\}_\epsilon=\epsilon^{-1}(f\star g-g\star f).
\]
The numerator is divisible by $\epsilon$, so this bracket
is integral. Write $\{f,g\}_0=f_qg_p-f_pg_q$ and
$X_f=\{f,-\}_0$. The quantum derivation
$\rho_f=\epsilon^{-1}\operatorname{ad}_\star f$ acts in
every chain slot, giving an operator $L_f^C$.

There is also a degree-minus-one insertion
\[
 h_f(a_0[a_1|\cdots|a_n])
 =-\sum_{i=0}^n(-1)^i
       a_0[a_1|\cdots|a_i|f|a_{i+1}|\cdots|a_n].
\]
On forms put $m_f=-df\wedge(-)$, also of degree minus one.
Constants vanish in both operations.
The brackets below are graded commutators of operators.

\begin{lemma}\label{mht:lem:cyclic-cartan}
The source and target operations satisfy
\begin{align*}
 [D_C,h_f]&=\epsilon L_f^C,& [B,h_f]&=0,&
 [L_f^C,h_g]&=h_{\{f,g\}_\epsilon},\\
 [D_\Omega,m_f]&=\epsilon L_{X_f},& [d,m_f]&=0,&
 [L_{X_f},m_g]&=m_{\{f,g\}_0}.
\end{align*}
Moreover $[h_f,h_g]=[m_f,m_g]=0$.
\end{lemma}
\begin{proof}
In $bh_f+h_fb$, every merger disjoint from the inserted $f$
occurs twice with opposite signs. The two neighboring mergers
leave the commutator with $f$ in each original slot.
Their sum is $\epsilon L_f^C$.
In $Bh_f+h_fB$, the inserted slot may occur on either side
of the cyclic origin. These terms cancel in pairs.
Inserting $f$ and $g$ in either order gives the same chain
with opposite signs, proving $[h_f,h_g]=0$.

A derivation applied to the insertion cancels against its
application to the original slots. The remaining term inserts
$\rho_f(g)$, giving the third source identity.
For forms, substitute the coordinate formulas for $\delta$.
They give $[\delta,-df\wedge]=L_{X_f}$.
Exterior differentiation and two odd exterior multiplications
anticommute, while Lie differentiation sends $dg$ to
$d\{f,g\}_0$. These prove the target identities.
\end{proof}

The classical vector fields do not represent the quantum
bracket strictly. Define
\[
 R(f,g)=\frac{\{f,g\}_\epsilon-\{f,g\}_0}{\epsilon},
 \qquad F_1(f)=L_{X_f},\qquad F_2(f,g)=dR(f,g)\wedge(-).
\]
The difference of brackets starts at order $\epsilon^2$,
so $R$ and $F_2$ are integral. Lemma~\ref{mht:lem:cyclic-cartan}
gives
\begin{equation}\label{mht:eq:cyclic-action-defect}
 [D_\Omega,F_2(f,g)]
   =[F_1(f),F_1(g)]-F_1(\{f,g\}_\epsilon).
\end{equation}

The ternary expression is the exterior differential of
\[
 \sum_{\mathrm{cyc}(f,g,k)}
       \bigl(\{f,R(g,k)\}_0+R(f,\{g,k\}_\epsilon)\bigr).
\]
Multiplying the sum by $\epsilon$ gives the difference of
the quantum and classical Jacobi identities. It is zero.
There are no further terms, since the odd wedge operators
$F_2$ anticommute.

A nonconstant correction makes the binary operation visible.
For $f=q^3/3$ and $g=p^4/4$,
\[
 \{f,g\}_\epsilon=q^2p^3+\tfrac12\epsilon^2p,
 \qquad F_2(f,g)=\tfrac12\epsilon\,dp\wedge(-).
\]
For the two cubics $q^3/3,p^3/3$, the bracket correction is
constant and $F_2$ vanishes. Their comparison homotopy below
still has a nonzero secondary value.

\section{Two integral comparison homotopies}

The first comparison homotopy has degree minus one:
\begin{equation}\label{mht:eq:cyclic-first-homotopy}
 \mathcal K_f
       =\epsilon^{-1}(\mathcal U h_f-m_f\mathcal U).
\end{equation}
At $\epsilon=0$, insertion under the normalized
Hochschild--Kostant--Rosenberg map is multiplication by $-df$.
Indeed its $n+1$ insertion positions cancel the extra factor
in $(n+1)!$. Hence the numerator is divisible by $\epsilon$.
The mixed and Cartan identities give
\[
 D_\Omega\mathcal K_f+\mathcal K_fD_C
     =\mathcal U L_f^C-L_{X_f}\mathcal U,
 \qquad d\mathcal K_f+\mathcal K_fB=0.
\]

The second comparison homotopy has degree minus two:
\begin{equation}\label{mht:eq:cyclic-second-homotopy}
 \mathcal K^{(2)}_{f,g}
 =\epsilon^{-2}
       \bigl(m_f\mathcal U h_g-m_g\mathcal U h_f
               -m_fm_g\mathcal U-\mathcal U h_fh_g\bigr).
\end{equation}
Its numerator vanishes on positive-length chains by exterior
degree. On a zero chain $a$ the numerator is
\begin{align*}
 &df\wedge\mathcal U_1(a[g])
   -dg\wedge\mathcal U_1(a[f])-a\,df\wedge dg\\
 &\hspace{35mm}-\mathcal U_2(a[f|g]-a[g|f]).
\end{align*}
Its constant coefficient is zero.

The one-bar map is even
in $\epsilon$. The linear coefficient of the two-bar
integral operator is
\[
 (P_{01}-P_{02}+P_{12})/12.
\]
Interchanging the last two slots reverses both this operator
and their exterior product. Antisymmetrization therefore
kills the linear coefficient. Thus
$\mathcal K^{(2)}$ is integral as well.

We specify the higher-action signs by the completed
Chevalley complex. Choose an enumeration $(f_a)$ of the
nonconstant monomials, viewed in $\mathfrak h_\epsilon$,
and let $c^a$ be the corresponding odd dual coordinates
of degree one. Define $f^a_{bc}\in R_\epsilon$ by
\[
 \{f_b,f_c\}_\epsilon=\sum_a f^a_{bc}f_a
 \quad\text{in }\mathfrak h_\epsilon.
\]

The scalar Chevalley differential is
\[
 d_{\mathrm{CE}}c^a=-\frac12\sum_{b,c}f^a_{bc}c^bc^c.
\]
An odd operator anticommutes with an odd variable from the
coefficient factor. The coordinate sums denote continuous
multilinear maps, coefficientwise in the parameters and in
fixed finite arity.

No finite basis is assumed. We use
product completion over arities.

\begin{theorem}\label{mht:thm:cyclic-hamiltonian-comparison}
The operators $F_1,F_2$ and zero higher components define an
$L_\infty$ action of $\mathfrak h_\epsilon$ on the mixed
form complex. In the stated Chevalley convention,
\[
 \mathbb D_\Omega
   =D_\Omega+d_{\mathrm{CE}}
       +\sum_a c^a F_1(f_a)
       -\sum_{a<b}c^ac^bF_2(f_a,f_b).
\]
The comparison
\[
 \mathbb U=\mathcal U-\sum_a c^a\mathcal K_{f_a}
                    +\sum_{a<b}c^ac^b
                              \mathcal K^{(2)}_{f_a,f_b}
\]
is an integral $L_\infty$ module quasi-isomorphism from the
chain action. It preserves the circle operators.
\end{theorem}
\begin{proof}
Temporarily invert $\epsilon$ and set
\[
 T=\epsilon^{-1}\sum_a c^a h_{f_a},\qquad
 M=\epsilon^{-1}\sum_a c^a m_{f_a}.
\]
These have degree zero and raise Chevalley arity.
Their exponentials therefore exist in the arity completion.
The Cartan identities give
\begin{align*}
 e^T(D_C+d_{\mathrm{CE}})e^{-T}
   &=D_C+d_{\mathrm{CE}}+\sum_a c^aL_{f_a}^C,\\
 e^M(D_\Omega+d_{\mathrm{CE}})e^{-M}
   &=\mathbb D_\Omega.
\end{align*}

On chains, the commutator with $d_{\mathrm{CE}}$ cancels
the quadratic term because the action and insertion use
the same bracket. On forms their discrepancy is $F_2$,
with the sign displayed in the theorem.
Further commutators vanish by anticommutation of the insertions.

Conjugate the mixed comparison:
\[
 e^M\mathcal Ue^{-T}=\mathbb U.
\]
Expansion through arity two gives exactly
\eqref{mht:eq:cyclic-first-homotopy} and
\eqref{mht:eq:cyclic-second-homotopy}, with the tensor signs
just specified. A term of arity $j$ sends a chain of length
$n$ to a form of degree $n+j$. Hence every term of arity
at least three vanishes on a surface.

Conjugation gives $\mathbb D_\Omega^2=0$ and
$\mathbb D_\Omega\mathbb U=\mathbb U\mathbb D_C$,
where $\mathbb D_C$ is the first conjugated differential.
All displayed components are integral, and their coefficient
modules have no $\epsilon$-torsion. The identities therefore
hold before localization. Their parts involving $u$ express
compatibility with the circle differentials.

The arity-zero comparison is a quasi-isomorphism by
Theorem~\ref{mht:thm:integral-cyclic-comparison}.
Its continuous cone contraction also contracts each fixed-arity
space of continuous multilinear maps by postcomposition.
The remaining differential raises arity. Perturbing in arity
therefore contracts the full cone in the product completion.
\end{proof}

The cubic first homotopy on a zero chain is
\[
 \mathcal K_{q^3/3}(a)=-\frac{\epsilon}{12}a_{pp}\,dq.
\]
The second homotopy has an all-order vacuum formula:
\begin{equation}\label{mht:eq:cyclic-secondary-vacuum}
 \mathcal K^{(2)}_{f,g}(1)
 =-\frac{\{f,g\}_\epsilon-\{f,g\}_0}{\epsilon^2}\Omega.
\end{equation}

To see this, the unit in the outer slot kills $P_{01}$
and $P_{02}$. Put $s=t_2-t_1$. The even moments of the
remaining operator are
\[
 2\int_0^1(1-s)(\tfrac12-s)^{2j}\,ds
       =\frac1{2^{2j}(2j+1)}.
\]
Together with $(2j)!$ and the two exterior derivatives,
these are the coefficients of the divided Moyal commutator.
The remaining terms of \eqref{mht:eq:cyclic-second-homotopy}
subtract its classical coefficient.

In particular,
\[
 \mathcal K^{(2)}_{q^3/3,p^3/3}(1)=-\Omega/6,
 \qquad
 \mathcal K^{(2)}_{q^3/3,p^3/3}(qp)=-qp\Omega/6.
\]
The second identity follows by the same terminating
simplex integral with outer coefficient $qp$.
These values persist at $\epsilon=0$ after the integral
division has been made. A strict classical projection by
itself does not retain this secondary comparison.

\section{The integrated Hamiltonian kernel}

The differential formulas also apply coefficientwise to smooth
inputs. On an oriented real plane use the volume $\Omega$.
For complex coordinates, include a fixed complementary
antiholomorphic volume form in each integral. Assume compact
support or periodicity, so integration by parts has no boundary
term. These are formal identities in $\epsilon$.

For scalar $f$ and matrix-valued inputs $a_0,a_1$, one has
\begin{equation}\label{mht:eq:integrated-cyclic-kernel}
 \int df\wedge\operatorname{Tr}\mathcal U_1(a_0[a_1])
     =\int\Omega\,\operatorname{Tr}(a_0\rho_f(a_1)).
\end{equation}
The remaining form factors and orientations agree on the two
sides.

To prove the identity, give $f,a_0,a_1$ momenta $k,l,m$.
Integration by parts imposes $k+l+m=0$.
If $c=\pi(k,m)$, then $\pi(l,m)=-c$, and the left multiplier is
\[
 c\frac{2\sinh(-\epsilon c/2)}{-\epsilon c}
     =\frac{2\sinh(\epsilon c/2)}{\epsilon}.
\]
This is the multiplier of $\rho_f$ on the right.
The Taylor coefficients prove the equality, with matrix
factors kept in their trace order.

The comparison for matrix algebras first uses the Morita trace
\[
 a_0[a_1|\cdots|a_n]\longmapsto
 \sum_{i_0,\ldots,i_n}
 (a_0)_{i_0i_1}[(a_1)_{i_1i_2}|\cdots|(a_n)_{i_ni_0}].
\]
Each adjacent matrix product sums the same repeated index on
both sides. Cyclic rotation preserves the closed index cycle.
The map therefore commutes with $b$ and $B$, and it descends
to normalized chains. In degree zero it sends $I_N$ to $N$,
retaining the rank normalization.

The integrated Hamiltonian insertion is a strictly cyclic
cochain. Put
\[
 \tau(a)=\int\Omega\,\operatorname{Tr}a,
 \qquad \phi_f(a_0,a_1)=\tau(a_0\star\rho_f(a_1)).
\]

Here the unit is allowed as a multiplier of compactly
supported inputs. Integration by parts gives
$\tau(a\star b)=\tau(ab)=\tau(b\star a)$ and
$\tau\rho_f=0$. Since $\rho_f$ is a derivation,
\[
 \phi_f(a_0,a_1)+\phi_f(a_1,a_0)
     =\tau\rho_f(a_0\star a_1)=0.
\]
The same derivation identity cancels its Hochschild boundary.
The equality $\rho_f(1)=0$ gives normalization.

Equation~\eqref{mht:eq:integrated-cyclic-kernel} expresses
this cochain as a form pairing. For additional graded
coefficients, the trace must be closed and $\rho_f$ must
commute with their differential. The graded identities then
have the corresponding Koszul signs.

\section{Continuous duals and cyclicity}

A general form current has no automatic cyclicity property.
We make the dual convention explicit. Let $C^n(A_\epsilon)$
be the coefficientwise continuous $R_\epsilon$-multilinear
functionals on $A_\epsilon^{n+1}$ which vanish when a bar
input is $1$. Define the scalar transposes by
\[
 (b\phi)(c)=\phi(bc),\qquad (B\phi)(c)=\phi(Bc).
\]
They have degrees $1,-1$ and satisfy the mixed identities.

In particular,
\begin{align*}
 (b\phi)(a_0,\ldots,a_{n+1})
   ={}&\sum_{i=0}^{n}(-1)^i
       \phi(a_0,\ldots,a_i\star a_{i+1},\ldots,a_{n+1})\\
    &+(-1)^{n+1}\phi(a_{n+1}\star a_0,a_1,\ldots,a_n).
\end{align*}
For $n\geq1$ the other operator is
\[
 (B\phi)(a_0,\ldots,a_{n-1})
   =\sum_{j=0}^{n-1}(-1)^{(n-1)j}
       \phi(1,a_j,\ldots,a_{n-1},a_0,\ldots,a_{j-1}).
\]
Set $B=0$ on $C^0$.
The completed scalar mixed complex is
$(C^\bullet(A_\epsilon)[[u]],b+uB)$, with $|u|=2$.
These transpose conventions fix the signs in its primitives.

A continuous functional $\beta$ on forms gives
\[
 \Phi_\beta(a_0,\ldots,a_n)
   =\langle\beta,\mathcal U_n(a_0[a_1|\cdots|a_n])\rangle.
\]
Transpose the two form operators by the same convention.
The mixed identities for $\mathcal U$ give a map from this
dual form complex to the scalar cochain complex.
A closed dual element gives a total mixed cocycle.
This assertion does not require products of distributions:
each evaluation has smooth or formal inputs and one
continuous functional.

A scalar one-cochain is \emph{strictly cyclic} when
$\Phi(a,b)=-\Phi(b,a)$. At $\epsilon=0$, the point current
\[
 \beta(a\,dq+b\,dp)=b(0)
\]
gives $\Phi_\beta(1,p)=1$ and $\Phi_\beta(p,1)=0$.
Thus an arbitrary continuous current fails this condition.
Smooth inputs equal to $1,p$ near the point give the same
example with compact support.

\begin{proposition}\label{mht:prop:all-order-current-cyclicity}
For a coefficientwise continuous functional
$\beta:\Omega_A^1[[\epsilon]]\to R_\epsilon$, the cochain
$\Phi_\beta(a,b)=\beta(\mathcal U_1(a[b]))$ is strictly
cyclic if and only if $\beta d=0$.
\end{proposition}
\begin{proof}
The series $S(P)$ in \eqref{mht:eq:circle-onebar} is even.
Interchanging the two coefficient slots therefore gives
\[
 \mathcal U_1(a[b])+\mathcal U_1(b[a])
       =d\bigl(mS(P)(a\otimes b)\bigr).
\]
Annihilation of exact forms proves sufficiency at every
parameter order. Conversely,
$\mathcal U_1(1[b])=db$ and $\mathcal U_1(b[1])=0$.
Cyclicity on this pair forces $\beta(db)=0$ for every $b$.
\end{proof}

\section{Point currents and two primitives}

Let $\mathscr J'_0$ be the vector space of finite linear
combinations of derivatives at $(q,p)=(0,0)$.
A series in $\mathscr J'_0[[\epsilon]]$ may have unbounded
derivative order as the parameter order increases.
For $\gamma$ in this space define
\[
 \beta_\gamma=\gamma\delta,\qquad
 \Phi_\gamma(a,b)=\gamma\delta\mathcal U_1(a[b]).
\]
This sends point functionals in degree zero to scalar
one-cochains in degree one. It specifies no output-valued
operation or multiplication of currents.

\begin{lemma}\label{mht:lem:point-current-classification}
Every coefficientwise finite point functional on one-forms
which annihilates exact forms is uniquely $\beta_\gamma$.
\end{lemma}
\begin{proof}
At one parameter coefficient write
\[
 \beta(a\,dq+b\,dp)
   =A(\partial_q,\partial_p)a(0)
       +B(\partial_q,\partial_p)b(0)
\]
for polynomials $A(X,Y),B(X,Y)$.
The condition $\beta d=0$ is the polynomial identity
$XA+YB=0$. Since $X,Y$ are relatively prime, there is a
unique polynomial $C$ with
\[
 A=-YC,\qquad B=XC.
\]

Take $\gamma(f)=C(\partial_q,\partial_p)f(0)$ and apply
the argument at every parameter coefficient.
For smooth complex-coordinate currents, antiholomorphic
derivative symbols may be included as polynomial coefficients.
The same regular-sequence argument still applies.
\end{proof}

The two-bar comparison supplies a second scalar cochain.
Let $\lambda_\gamma(c\Omega)=\gamma(c)$ and put
\[
 \Psi_\gamma(a_0,a_1,a_2)
      =\lambda_\gamma\mathcal U_2(a_0[a_1|a_2]).
\]
It is a normalized cochain of degree two.
The integral mixed map relates its Connes boundary to
$\Phi_\gamma$.

\begin{theorem}\label{mht:thm:point-current-two-primitives}
The point cochains satisfy
\[
 b\Phi_\gamma=B\Phi_\gamma=0,\qquad
 \epsilon\Phi_\gamma=b\gamma,\qquad
 b\Psi_\gamma=0,\qquad B\Psi_\gamma=\Phi_\gamma.
\]
Consequently
\begin{equation}\label{mht:eq:point-current-primitives}
 (b+uB)\gamma=\epsilon\Phi_\gamma,
 \qquad (b+uB)\Psi_\gamma=u\Phi_\gamma.
\end{equation}
The class $[\Phi_\gamma]$ is nonzero exactly when the
constant $\epsilon$ coefficient of $\gamma$ is nonzero.
The resulting copy of $\mathscr J'_0$ in degree-one mixed
cohomology is annihilated by both $\epsilon$ and $u$.
\end{theorem}
\begin{proof}
Direct differentiation of the one-bar formula gives
\[
 \delta\mathcal U_1(a[b])
       =mP S(P)(a\otimes b)
       =\frac{a\star b-b\star a}{\epsilon}.
\]
This proves $\epsilon\Phi_\gamma=b\gamma$.
Associativity gives $b^2\gamma=0$, and the scalar coefficient
module has no $\epsilon$-torsion. Hence $b\Phi_\gamma=0$.
Both unit evaluations vanish, so
$B\Phi_\gamma(a)=\Phi_\gamma(1,a)=0$.

On a chain of length three,
$\mathcal U_2 b=\epsilon\delta\mathcal U_3=0$.
Transposing proves $b\Psi_\gamma=0$.
On a chain $a[b]$, the second mixed identity gives
\[
 (B\Psi_\gamma)(a,b)
   =\lambda_\gamma d\mathcal U_1(a[b])
   =\gamma\delta\mathcal U_1(a[b])
   =\Phi_\gamma(a,b).
\]
The middle equality uses
$d(a\,dq+b\,dp)=(b_q-a_p)\Omega$.
This proves the two independent primitives in
\eqref{mht:eq:point-current-primitives}.

In total degree one a boundary can only come from a scalar
zero-cochain, since cochain degrees and powers of $u$ are
nonnegative. Modulo $\epsilon$, its Hochschild boundary
vanishes because $A$ is commutative.
If the constant coefficient $\gamma_0$ is nonzero, choose
a polynomial $f$ with $\gamma_0(f)\ne0$ and a polynomial
$g$ with $\partial_pg=f$.
Then $\Phi_{\gamma_0}(q,g)=\gamma_0(f)\ne0$.
Thus $\Phi_\gamma$ cannot be a boundary.

Conversely, if $\gamma=\epsilon\gamma_1$, then
$\Phi_\gamma=b\gamma_1$ is a mixed boundary.
The assignment of a class therefore depends exactly on
$\gamma_0$ and embeds $\mathscr J'_0$ as stated.
The two displayed primitives prove its annihilation by
each parameter.
\end{proof}

For $\gamma=\operatorname{ev}_0$, the normalizations are
\[
 \Phi_\gamma(q,p)=1,\qquad
 \Phi_\gamma(q^3,p^3)=\tfrac32\epsilon^2,\qquad
 \Psi_\gamma(1,q,p)=\tfrac12.
\]
The first primitive uses a zero-cochain and multiplication
by $\epsilon$. The second uses a two-cochain and the actual
Connes operator. They have different degrees and remain
distinct constructions. Neither identifies this scalar
current with a geometric normal residue or a closed field.

\section{Transport by a quantum shear}

A nonlinear coordinate motion tests the comparison homotopy.
For $f=q^3/3$,
\[
 \rho_f=q^2\partial_p+\frac{\epsilon^2}{12}\partial_p^3.
\]
Its two terms commute. The corresponding quantum flow is
\[
 \Phi_t(a)=e^{t\epsilon^2\partial_p^3/12}a(q,p+tq^2).
\]
The exponential terminates on each polynomial input, and
formal inputs are treated coefficientwise in the specified
completion. Since $\rho_f$ is a derivation, $\Phi_t$ is a
Moyal automorphism.

Its value on a cubic is
\[
 \Phi_t(p^3)=(p+tq^2)^3+\tfrac12t\epsilon^2.
\]
Ordinary pullback alone would omit the last term.

Let $\phi_t^*$ be the classical pullback on forms, and
let $\Phi_t$ act in each chain slot. Define
\begin{equation}\label{mht:eq:cyclic-flow-homotopy}
 \mathcal H_t=\int_0^t\phi_{t-s}^*\mathcal K_f\Phi_s\,ds.
\end{equation}
For polynomial inputs, each parameter coefficient of the
integrand is polynomial in the flow parameter. For formal
inputs this holds after fixing the coordinate coefficient too.
Integration of these coefficients and the first homotopy
identity give
\[
 \mathcal U\Phi_t-\phi_t^*\mathcal U
       =D_\Omega\mathcal H_t+\mathcal H_tD_C.
\]

Indeed differentiating $\phi_{t-s}^*\mathcal U\Phi_s$
with respect to $s$ gives
$\phi_{t-s}^*(\mathcal U L_f^C-L_{X_f}\mathcal U)\Phi_s$.
The mixed differentials commute with both flows.
Integration supplies the two endpoints and proves the formula.
For several Hamiltonian motions the second comparison
homotopy remains part of their compatibility data.

\section{A logarithmic torus}

The same kernel has a global example without choosing
logarithms. Let $A=\mathbb C[x^{\pm1},y^{\pm1}]$, and put
\[
 D_x=x\partial_x,\qquad D_y=y\partial_y,\qquad
 \alpha=d\log x,\qquad\beta=d\log y.
\]
Define its formal product on Laurent monomials by
\[
 x^ay^b\star x^cy^d
    =e^{\epsilon(ad-bc)/2}x^{a+c}y^{b+d}.
\]
Bilinearity of the exponent proves associativity.

The constant Laurent coefficient is a cyclic trace: only
opposite modes contribute, and their alternating pairing is
zero. Replace $\partial_q,\partial_p$ by $D_x,D_y$ in
\eqref{mht:eq:circle-comparison}.
The commuting derivations give the same mixed identities
globally, without logarithmic functions in the algebra.

The undeformed diagonal Koszul resolution of this Laurent
algebra is the localization of the polynomial diagonal
resolution. It proves the Hochschild comparison, and the
same completed filtration argument proves the mixed
quasi-isomorphism. Each parameter coefficient has finite
Laurent support. The exponential is now a formal parameter
series, rather than a terminating polynomial expression.

After inverting $\epsilon$ and then completing in $u$,
consider the Laurent mode $(m,n)\in\mathbb Z^2$.
The form differential on this mode is
\[
 D_{m,n}=\epsilon\iota_{-nD_x+mD_y}
                    +u(m\alpha+n\beta)\wedge.
\]
For a nonzero mode put
\[
 H_{m,n}=\frac{-n\alpha+m\beta}{\epsilon(m^2+n^2)}\wedge.
\]

Contraction with $-nD_x+mD_y$ evaluates the numerator to
$m^2+n^2$. The two exterior multiplication operators
anticommute. Therefore
\[
 D_{m,n}H_{m,n}+H_{m,n}D_{m,n}=1.
\]
The zero mode has zero differential. The mixed cohomology is
consequently
\begin{equation}\label{mht:eq:torus-cyclic-cohomology}
 \mathbb C((\epsilon))[[u]]\otimes\Lambda(\alpha,\beta).
\end{equation}
The homotopy preserves Laurent support and lowers the
$\epsilon$ bound by one, so it is defined on the stated
coefficientwise completion.

The derivations $D_x,D_y$ are outer on the Laurent algebra.
An inner derivation preserving the mode of each input could
only use a zero-mode Hamiltonian. Distinct nonzero Hamiltonian
modes shift each input to distinct output modes. A zero-mode Hamiltonian is
constant and has zero commutator. Thus local Hamiltonian
coordinates do not remove these global vector fields or
the four zero-mode form classes.

\section{The remaining source comparison}

The circle map and its Hamiltonian homotopies are constructions
for the associative coefficient algebra. The integrated
Hamiltonian insertion gives a closed cyclic cochain under its
trace and support hypotheses. A general distribution supplies
a dual mixed cochain, with strict cyclicity governed by
Proposition~\ref{mht:prop:all-order-current-cyclicity}.
The point-current primitives are scalar operations on this
specified recipient.

To obtain the mixed field-theory comparison, one must supply
a map from the geometric closed-source complex. It must
preserve its differential, interaction, boundary restriction
and polar collision operations. The scalar parameter
primitives above neither choose this map nor fix its physical
normalization. Tensoring with a normal complex preserves the
mixed identities, but cyclic insertions still require a
compatible normal pairing and the stated tangential cyclicity.

For a general holomorphic symplectic surface, coordinate
changes must carry the comparison and its higher homotopies
through overlaps. A connection, deformation period and trace
normalization enter this descent. The explicit shear supplies
one transport homotopy, while the logarithmic torus exhibits
global classes which a local calculation must retain.
Neither example constructs the full descent or its supported
chiral comparison.
